\section{Alcance y delimitación}
~\label{sec:alcance}

El alcance del proyecto se define sobre segmentos de red 
de pequeña o mediana escala administrados por un único 
responsable, que accede al sistema desde una aplicación móvil y cuenta 
con conocimientos técnicos en administración de redes. Los sensores de 
captura se despliegan en los endpoints bajo control del administrador, 
mientras que los dispositivos ajenos a esa gestión solo son objeto de 
detección e identificación básica a partir de técnicas de 
descubrimiento activo.

\subsection{Inclusiones}

El proyecto comprende el desarrollo de los cuatro componentes 
establecidos en el acta de constitución:

\begin{itemize}
    \item \textbf{Sensores de captura desplegables en endpoints}, 
    responsables de capturar el tráfico de red de la interfaz en la 
    que están instalados y transmitirlo hacia un servidor colector 
    centralizado dentro de la misma red local, recolectar métricas 
    básicas del estado del endpoint como uso de CPU y memoria, y 
    ejecutar descubrimiento básico de otros dispositivos presentes en 
    el segmento de red mediante técnicas como ARP.@

    \item \textbf{Backend de análisis centralizado}, encargado de 
    recibir y almacenar los datos enviados por los sensores, procesar 
    el tráfico capturado para generar métricas contextualizadas como 
    consumo de ancho de banda, retransmisiones TCP, tiempos de 
    respuesta DNS y perfiles de tráfico por dispositivo, construir y 
    mantener el inventario dinámico de activos, identificar 
    comportamientos inusuales y exponer la información procesada a 
    través de una API consumida por la aplicación móvil.

    \item \textbf{Aplicación móvil de consulta y notificación}, que 
    permite al administrador consultar el inventario de dispositivos, 
    visualizar métricas tanto a nivel de segmento de red como por 
    endpoint individual, contextualizar los dispositivos dentro del 
    segmento como representación de la topología, y recibir 
    notificaciones push ante comportamientos de tráfico inusuales.

    \item \textbf{Infraestructura de integración y entrega continua}, 
    que gestiona la compilación, pruebas automatizadas y despliegue 
    del backend, así como la publicación de la aplicación móvil.
\end{itemize}

\subsection{Exclusiones}

El proyecto no contempla los siguientes aspectos, definidos 
explícitamente como fuera de alcance:

\begin{itemize}
    \item Mecanismos de respuesta automática ante incidentes; el 
    sistema monitorea y notifica, pero no interviene sobre la red.

    \item Administración activa de la red, como denegar o permitir 
    acceso a dispositivos.

    \item Gestión o control remoto de los dispositivos conectados.

    \item Monitoreo del tráfico de dispositivos que no cuenten con 
    sensor instalado; para estos, el sistema se limita a su detección 
    e identificación básica dentro del inventario.

    \item Infraestructuras de gran escala, como redes con múltiples 
    sedes, segmentos WAN o despliegues en nubes públicas.

    \item Almacenamiento histórico ilimitado de datos de tráfico; el 
    sistema opera con ventanas de retención acotadas por la capacidad 
    de la infraestructura disponible.
\end{itemize}